% ============================================================================
% Fichier cible : ARIEL_LOGMINER_MEMOIRE_FINAL/chapters/chapitre2_etat_art.tex
% Insertion : immédiatement avant « \section{Systèmes multi-agents} ».
% Agent : Agent 3 — État de l'art et positionnement, chantier 2026-11-05.
% Sources : documentations officielles Wazuh, Elastic Security et OpenSearch Security
%   Analytics ; article LogAI ; article et dépôt Loglizer.
% ============================================================================

\section{Cadres intégrés et bibliothèques d'analyse de journaux}

L'intégration de l'ingestion, de la détection, de la corrélation et de la restitution
n'est pas une lacune générale de l'offre existante. Plusieurs plateformes réunissent déjà
ces fonctions, tandis que des bibliothèques de recherche unifient le parsing et
l'évaluation de détecteurs. Elles constituent des références de positionnement, même si
leurs objectifs, leurs coûts et leurs protocoles ne sont pas identiques à ceux d'un
prototype académique léger.

Wazuh associe des agents de collecte à un serveur d'analyse, un indexeur et un tableau de
bord ; son architecture documente aussi des déploiements centralisés ou en grappe
\cite{wazuh2026architecture}. Elastic Security s'appuie sur des données normalisées avec
ECS et propose des règles de détection, de la corrélation d'événements, des alertes et des
outils d'investigation \cite{elastic2026ecs,elastic2026security}. OpenSearch Security
Analytics organise les journaux par types, applique des règles aux détecteurs, produit
des constats et des alertes, et inclut un moteur de corrélation entre types de journaux
\cite{opensearch2026securityanalytics}. Ces trois cadres couvrent donc déjà une part
substantielle de la chaîne fonctionnelle revendiquée par Ariel Logminer.

Du côté de la recherche, Loglizer fournit une boîte à outils de détection d'anomalies
fondée sur l'apprentissage automatique et associée à une comparaison de méthodes sur des
journaux système \cite{he2016experience,logpai2026loglizer}. LogAI propose une interface
commune pour le résumé, le regroupement et la détection d'anomalies, plusieurs familles de
modèles, le modèle de données OpenTelemetry et une interface interactive
\cite{cheng2023logai}. Ces bibliothèques réduisent le travail nécessaire pour comparer des
algorithmes, mais ne se présentent pas comme des SIEM complets.

\begin{table}[H]
\centering
\caption{Cadres intégrés et bibliothèques proches du périmètre étudié}
\label{tab:cadres-integres}
\scriptsize
\begin{tabularx}{\textwidth}{p{2.25cm}p{2.1cm}p{3.5cm}X}
\toprule
\textbf{Cadre} & \textbf{Nature} & \textbf{Fonctions documentées pertinentes} &
\textbf{Différence avec l'objet de ce mémoire} \\
\midrule
Wazuh \cite{wazuh2026architecture} & Plateforme de sécurité &
Collecte par agents, analyse, indexation, alertes et tableau de bord ; déploiement en grappe. &
Plateforme opérationnelle ; ce mémoire ne la compare ni en débit ni en qualité de détection. \\
Elastic Security \cite{elastic2026security} & SIEM et sécurité des terminaux &
Normalisation ECS, règles, corrélation, alertes, investigation et réponse. &
Écosystème intégré et étendu ; Ariel Logminer vise une instrumentation expérimentale plus réduite. \\
OpenSearch Security Analytics \cite{opensearch2026securityanalytics} & Extension SIEM d'OpenSearch &
Types de journaux, règles, détecteurs, constats, alertes et corrélation. &
Alternative directe pour plusieurs fonctions ; aucune comparaison expérimentale n'est réalisée ici. \\
LogAI \cite{cheng2023logai} & Bibliothèque de recherche &
Prétraitement, résumé, regroupement, détection, modèles classiques et profonds, interface. &
Pas d'allocation Contract Net ; cadre plus large pour comparer les méthodes d'analyse. \\
Loglizer \cite{he2016experience,logpai2026loglizer} & Boîte à outils de recherche &
Extraction de caractéristiques et détecteurs d'anomalies sur journaux. &
Centré sur l'analyse algorithmique, sans chaîne SOC intégrée. \\
\bottomrule
\end{tabularx}
\end{table}

Cette comparaison retire toute prétention à combler un vide fonctionnel. Ariel Logminer
ne remplace pas ces cadres et aucune expérience ne démontre qu'il les surpasse. Son objet
défendable est plus étroit : rendre explicites, dans une chaîne locale, les unités
d'analyse, les protocoles de séparation, le routage, l'allocation et les traces de preuve,
afin d'étudier les limites d'une intégration légère. L'originalité éventuelle doit être
cherchée dans cette instrumentation et dans les comparaisons contrôlées, pas dans la
simple coexistence des fonctions.

